\documentclass[12pt,letterpaper]{article}
\usepackage{listings}

\lstset{breaklines=true, columns=flexible, basicstyle=\small\ttfamily}

\begin{document}

\title{New and Improved Quick Start Guide for Revised Tfit}
\author{Michael Gohde}
\date{October 11, 2018}
\maketitle

\abstract{Tfit is a tool that attempts to find useful parameters in sequencing datasets.}

\section{Building Tfit-revisions}
In order to start working tiwh Tfit, it is first necessary to build it. Tfit depends on the following libraries to be built:

\begin{enumerate}
\item MPI (can be provided through OpenMPI)
\item GCC $>=$5.0
\item mpic++ (provided through OpenMPI)
\item (optional) docker
\end{enumerate}

In a shared environment such as Fiji, all that is necessary to build and run the latest version of Tfit is to:
\begin{enumerate}
\item chdir into the src directory
\item run `module load gcc/7.1.0'
\item run `module load openmpi'
\item run `make -j(some appropriate number of jobs)'
\end{enumerate}

On a personal machine, these requirements can be satisfied through appropriate package installations. All that is necessary in that case is to follow the first and last steps specified above.

\section{Invoking Tfit}
Tfit has several modes of operation referred to here as \textit{modules}. These modules perform different tasks, but many feature similar command line arguments. These overlaps will be noted as appropriate in this guide.

Invoking Tfit without any parameters will result in the following:
\begin{lstlisting}
Usage: Tfit modulename [arguments] [outputarguments]
Where modulename is one of the following:
        bidir - This module searches the genome for areas resembling bidirectional transcription
                by comparing a fixed template mixture model to a noise model by a log-likelihood ratio score.
        bidir_old - This module attempts to behave as documented in published versions of Tfit
        model - This module attempts to generate an optimal set of parameters per region instead 
                of using a fixed set of parameters for the entire genome


Where [arguments] is one or more of the following for the bidir or bidir_old modules:
        -i      Forward bedgraph file
        -j      Reverse bedgraph file.
        -ij     Joint (both forward and negative) bedgraph file. This parameter may be used in place of -i and -j if reads are in one bedgraph file.
                Please note that bedgraphs used by Tfit should be produced using "bedtools genomecoverage" or equivalent.

Additional (optional) parameters for the bidir module:
        -tss    Transcription model path. Models are provided for hg19 and mm10 in the annotations directory of this project.
        -chr    Run bidir only on the specified chromosome by name. The default is, "all"
        -fdr    Generate a likelihood score distribution on the input data. This parameter has yet to be fully documented and tested.
        -mle [0, 1]     Automatically invoke the model module after generating preliminary bidirectional predictions.


Where [arguments] is one or more of the following for the model module:
        -i      Forward bedgraph file
        -j      Reverse bedgraph file
        -ij     Both forward and reverse bedgraph file. This parameter may be used in place of -i and -j if reads are in one bedgraph file.
        -k      Bedgraph file containing a set of regions of interest.


Where [outputarguments] is ALL of the following:
        --prelimhits, -ph       Output file for preliminary filtering in the bidir module.
        --models, -m    Output file for models computed in the model module OR bidir module with the -mle parameter set to 1.
        --bidirpreds,-bp        Bidirectional predictions file (used to represent "regions of interest" for the model module.
        --logfile,-lf   Log file in which to write error and message output.


(To get a full usage statement, run TFit with the -h parameter.)
\end{lstlisting}

The above listing is Tfit's \textit{usage statement}. This presents a quick overview of all of Tfit's modules and their parameters. From it, we can determine that there are presently three available modules:
\begin{enumerate}
\item bidir -- This module attempts to fit a distribution to existing data.
\item bidir\_old -- This module is the same as bidir, except it attempts to behave more similarly to early versions of Tfit. User testing has shown that this module generally tends to produce a slightly higher number of overall calls than `bidir.'
\item model -- This module attempts to model every given bidirectional in a dataset individually instead of applying precomputed parameters.
\end{enumerate}

From the usage statement, we can begin to see how Tfit is to be invoked. It is first necessary to call Tfit, then to specify the name of a module. After this, various arguments specific to the module followed by arguments related to how output is to be handled are to be specified. 

\subsection{Invoking the Bidir (and Bidir\_Old) Module(s)}
\textit{Please note that bidir\_old behaves similarly enough to bidir that it is invoked in exactly the same way.}

The bidir module requires, at minimum, the following information to run:
\begin{enumerate}
\item An input bedgraph (this may be two bedgraphs if each covers only one strand).
\end{enumerate}

Input bedgraphs are specified with the \textit{-i, -j,} and \textit{-ij} parameters. \textit{-i} specifies a bedgraph with only forward strand data, while \textit{-j} specifies a bedgraph with only reverse strand data. \textit{-i}j can be used in place of \textit{-i} or \textit{-j}, as it specifies a bedgraph file with both forward and reverse strand data.

While it is possible to run Tfit with only the minimum number of parameters (technically just a bedgraph file specification), it may produce unexpected, difficult to find, or incorrect results. As such, it is advised to further specify the following:

\begin{enumerate}
\item A complete set of output filenames. (These will be detailed following this)
\item A transcription model. Tfit uses this file to infer parameters about a given dataset before selecting various regions of interest. Models are provided for hg19 and mm10 in the annotations directory of the Tfit repository. This is specified with the `-tss' parameter.
\item A chromosome name. By default, Tfit will attempt to fit every chromosome in a given dataset, which tends to be extremely computationally intensive. This parameter matches a single chromosome specification in a given input bedgraph (eg. `chr1' etc.) This is specified with the `-chr' parameter.
\item Whether to run the model module on bidir's inferences. This removes the need to call fstitch or a related program to produce a set of preliminary reads for the model module, as it is possible to have Tfit automatically run the model module after inferring information about a given dataset. This feature tends to be very useful when trying to find regions of transcription as quickly as possible. If this is to be enabled, then Tfit must be invoked with `-mle 1'.
\end{enumerate}

Originally, Tfit would automatically generate output filenames given a job name and an output directory. This automatic naming convention unfortunately led to files that were difficult to distinguish from one another after running multiple jobs in sequence. As such, the present implementation requires that users manually specify the full names of all relevant output files. While some of the following may be omitted depending on the specific arguments and module specification desired, it should be considered good form to specity \textit{all} of the following for every Tfit run:

\begin{enumerate}
\item \textit{--prelimhits} - The bidir module produces a set of preliminary predictions as to which data points ``look'' like they might be useful areas to call. This bed file is intended to store all of those predictions.
\item \textit{--models} - When it fits each call, the model module generates a TSV (tab separated values) file containing a set of all models applied to the input dataset as well as their performance. 
\item \textit{--bidirpreds} - This represents the final set of model calls as stored in a bed file. Each call is labeled with a name and various model parameters. This file should be readable by genome browser utilities like IGV.
\item \textit{--logfile} - This represents a plain text file into which Tfit will write a transcript of the model fitting and bidirectional prediction process. 
\end{enumerate}

Please note that although output parameters are shown to exist \textit{after} all other parameters by the usage statement, they may be specified in any order following the module specification.


Given the above information, it is possible to construct an example Tfit invocation:
\begin{lstlisting}
Tfit bidir -ij /path/to/my.bed --prelimhits prelim.bed --models models.tsv --bidirpreds preds.bed --logfile log.log
     -tss hg19.tss -chr chr1 -mle 1
\end{lstlisting}

The above invocation will generate a set of preliminary bidirectional reads over chromosome 1 given the HG19 annotation file present in this repository. When it is finished generating and writing these bidirectionals, it will proceed to run the model module to fit each of the regions called.

\subsection{Invoking the Model Module}
The model module requires, at minimum, the following information to run:
\begin{enumerate}
\item An input bedgraph file.
\item A bedgraph file specifying regions of interest. This is specified with the `-k' parameter.
\end{enumerate}

While it is possible to run Tfit's model module with only the minimum parameters listed above, it may be helpful to specify some of the optional parameters listed below:

\begin{enumerate}
\item A complete set of output filenames. (These were detailed in the subsection on invoking the \textit{Bidir} module. Please note that the same process applies.)
\item A maximum number of iterations to compute. This is useful when the model module doesn't converge immediately, and it is specified with the `-mi' parameter.
\item A convergence threshold. This is useful when the model module doesn't converge or when a tighter tolerance is necessary. This is specified with the `-ct' parameter.
\end{enumerate}

There exist several other parameters, the behaviors of which have not been explored by the author of this document at present. Unfortunately, the bidir module has been the topic of significantly more development, though further documentation for the model module will be written later. 

The \textit{-k} parameter can also be used with the preliminary bidirectional outputs from the bidir module, though it is recommended to simply use the \textit{-mle 1} parameter instead. In order to quickly start with the bidir module, it should be possible to use the default tss annotation files provided with this project as regions of interest for the model module. 

Based on the above listings, it is possible to construct an example Tfit invocation:
\begin{lstlisting}
Tfit model -ij /path/to/my.bed -k /path/to/regions.bed 
     --prelimhits prelim.bed --models models.tsv --bidirpreds preds.bed --logfile log.log
\end{lstlisting}

\end{document}
